\section{Derivative energy and curvature concentration}
\label{sec:energy-curvature}

We retain the reduced, distinct-factor hypotheses of
Theorem~\ref{thm:actual-comparison}. All assertions below
concern the actual solution \(H(T)\).

\begin{theorem}\label{thm:energy-curvature}
For sufficiently large \(T\),
\begin{equation}\label{eq:actual-energy-bounds}
 \|\Phi_{H(T)}\|_\infty\le\frac{C}{1+T},\qquad
 \frac{c}{(1+T)^2}\le\mathcal E(T)\le\frac{C}{(1+T)^2}.
\end{equation}
At each \(p\in\Sigma\), the convergence
\[
 H(T,r_p(T)y)\longrightarrow H_{\mathrm{loc}}(y)
\]
holds in \(C^\infty\) on every compact \(y\)-set, in the
fixed holomorphic frames of Lemma~\ref{lem:normalized-reference}.
In particular, for every fixed \(R>0\),
\begin{equation}\label{eq:actual-core-energy}
 \int_{B_{Rr_p(T)}(p)}|F_{H(T)}|^2\,\dd V
 \longrightarrow8\pi^2\mathfrak F(R),\qquad
 \mathfrak F(R)=\frac{R^4(R^2+3)}{(1+R^2)^3}.
\end{equation}
The first radius \(r_{F,p}(T)\) where the curvature energy
in the coordinate ball about \(p\) reaches \(4\pi^2\)
satisfies
\begin{equation}\label{eq:half-charge-radius}
 r_{F,p}(T)=r_p(T)(1+o(1)).
\end{equation}
\end{theorem}

\subsection{The Donaldson functional and mean-curvature decay}

Use the Donaldson functional with an arbitrary additive constant
and the first variation
\begin{equation}\label{eq:donaldson-variation}
 \dd\mathcal D_H(Hs)=\int_X\tr(\Phi_Hs)\,\dd V,
 \qquad \dd V=\omega^2/2 .
\end{equation}
The first and second variations are given in
\cite[(6.3.27), (6.3.29)--(6.3.34)]{Kobayashi87}.
On a metric geodesic \(H_\theta=He^{\theta s}\), its
acceleration is zero, so the general second-variation
formula and integration by parts give
\begin{equation}\label{eq:donaldson-convexity}
 \frac{\dd^2}{\dd\theta^2}\mathcal D(H_\theta)
       =\|D'_{H_\theta}s\|_2^2\ge0 .
\end{equation}
No HE endpoint is needed for this calculation. The actual
flow in our time convention satisfies
\begin{equation}\label{eq:donaldson-dissipation}
 \frac{\dd}{\dd T}\mathcal D(H(T))
       =-2\|\Phi_{H(T)}\|_2^2.
\end{equation}
All three identities use the same volume normalization.

\begin{lemma}\label{lem:functional-decay}
The functional has a finite infimum \(\mathcal D_{\inf}\)
on the smooth metrics of determinant \(h_{\det}\), and
\begin{equation}\label{eq:functional-decay}
 0\le\mathcal D(H(T))-\mathcal D_{\inf}\le C/T
\end{equation}
for large \(T\).
\end{lemma}

\begin{proof}
The corrected family has
\(\|\Phi_{\widehat H_a}\|_\infty\le Ca\) and bounded
logarithmic-parameter velocity, by
\eqref{eq:corrected-bounds} and \eqref{eq:corrected-tangent}.
Thus \eqref{eq:donaldson-variation} gives
\(|\partial_a\mathcal D(\widehat H_a)|\le C\).
There is a finite limit \(\mathcal D_{\inf}\), with
\[
 |\mathcal D(\widehat H_a)-\mathcal D_{\inf}|\le Ca.
\]
For a fixed smooth determinant-fixed metric \(H\), the
supremum distance to \(\widehat H_a\) is at most
\(C_H+C|\log a|\). Indeed the reference eigenvalues
\(w,w^{-1}\) lie between fixed positive multiples
of \(a\) and \(a^{-1}\), and the correction is
\(O(a\ell(a))\). Convexity at \(\widehat H_a\) gives
\[
 \mathcal D(H)\ge\mathcal D(\widehat H_a)
 +\int_X\tr\bigl(\Phi_{\widehat H_a}
             \log(\widehat H_a^{-1}H)\bigr)\,\dd V .
\]
The final term is \(O_H(a\ell(a))\). Letting \(a\downarrow0\)
with \(H\) fixed proves that \(\mathcal D(H)\ge\mathcal D_{\inf}\).
The comparison metrics attain this value in the limit,
so it is the infimum.

For any fixed smooth metric \(K\), let
\(d_2(H,K)^2=\int_X|\log(H^{-1}K)|_H^2\,\dd V\).
The endpoint variation of squared distance and convexity give
\begin{equation}\label{eq:evi}
 \frac12\frac{\dd}{\dd t}d_2(H(t),K)^2
 \le2\{\mathcal D(K)-\mathcal D(H(t))\}.
\end{equation}
To check the sign and factor, the gradient of half squared
distance at \(H\) is minus the logarithm toward \(K\);
the velocity is \(-2\Phi_H\); and convexity bounds the
pairing of \(\Phi_H\) with that logarithm by
\(\mathcal D(K)-\mathcal D(H)\).

For each final \(T\), fix \(K=\widehat H_{a(T)}\) while
integrating \eqref{eq:evi} over \([T/2,T]\).
The initial distance in that interval is uniformly bounded:
Theorem~\ref{thm:actual-comparison} compares \(H(T/2)\) with
\(\widehat H_{a(T/2)}\), and the bounded parameter velocity
bounds the remaining distance by
\(C|\log(a(T/2)/a(T))|\). The ratio tends to two.
Since \(\mathcal D\) decreases along the actual flow,
integration gives
\(\mathcal D(H(T))-\mathcal D(K)\le C/T\).
Adding the nonnegative comparison gap, at most \(Ca(T)\),
proves \eqref{eq:functional-decay}.
\end{proof}

The mean-curvature identity \eqref{eq:mean-curvature-heat}
makes \(\|\Phi(T)\|_2^2\) nonincreasing. Therefore
\[
 \frac T2\|\Phi(T)\|_2^2
 \le\int_{T/2}^T\|\Phi(t)\|_2^2\,\dd t
 \le\tfrac12\{\mathcal D(H(T/2))-\mathcal D_{\inf}\}
 \le C/T .
\]
The unit-time scalar heat \(L^1\)-to-\(L^\infty\) estimate
applied to \(|\Phi|^2\) gives
\(\|\Phi(T)\|_\infty^2\le C\|\Phi(T-1)\|_2^2\).
This proves the first estimate of
\eqref{eq:actual-energy-bounds}.
The exact integrated inclusion identity then gives
\[
 \mathcal E(T)
 =\int_X\langle\Phi_H\iota,\iota\rangle\,\dd V
 \le\|\Phi_H\|_\infty Z(T)\le C/T^2 .
\]
The lower bound follows from \eqref{eq:reduced-liminf}
and \eqref{eq:actual-mass}. This argument estimates the
whole inclusion derivative and does not differentiate
the mass equivalent.

\subsection{Smooth convergence on the intrinsic scale}

Fix \(p\), a sequence \(T_n\to\infty\), and a finite
\(R>0\). Put \(r_n=r_p(T_n)\) and use the rescaled
cylinder \(B_{4R}(0)\times[-4,4]\), corresponding to
\[
 z=r_ny,\qquad T=T_n+r_n^2\tau .
\]
For large \(n\) it lies at positive times in the
original coordinate ball. The rescaled base forms
\(\omega_n=r_n^{-2}\delta_n^*\omega\), where
\(\delta_n(y)=r_ny\), converge smoothly to the
Euclidean form. Choose the time-independent scalar
reference matrix
\[
 H_{0,n}(y)=\sqrt{h_L(r_ny)h_M(r_ny)}\,I .
\]
It has the exact rescaled determinant, uniformly
bounded smooth coefficients and inverse, and zero
trace-free mean curvature.

Set \(G_n(y,\tau)=H(T_n+r_n^2\tau,r_ny)\).
These matrices converge uniformly on the cylinder
to \(H_{\mathrm{loc}}(y)\). Indeed \(a'=O(a^2)\)
implies \(a(T_n+r_n^2\tau)/a(T_n)\to1\) uniformly,
\(\alpha(T_n)/a(T_n)\to1\), and the global metric
comparison \eqref{eq:actual-static-comparison} tends
to equality throughout the indicated times.
The explicit reference metrics have the stated limit.
In particular \(G_n\) and \(H_{0,n}\) are uniformly
comparable, and the relative logarithm is trace-free.
The rescaled projective mean curvature is
\(r_n^2\Phi_{H(T_n+r_n^2\tau)}\). It is uniformly
bounded by \eqref{eq:actual-energy-bounds}; the original
maximum-principle bound would already suffice here.

Apply the unit-radius \(k=0\) interior estimate of
\cite[Appendix~C, Theorem~C.1 and Proposition~C.2]{JacobWalpuski18}
on nested balls at each time. Its inputs are bounded
base and reference geometry, a bounded trace-free
metric logarithm, and bounded projective mean curvature.
All hold uniformly in these rescaled charts. The
conclusion bounds the logarithm in \(W^{2,p}\) on
smaller balls for every finite \(p\), hence in
\(C^{1,\beta}\) for \(p>4\). A finite covering gives
a uniform spatial gradient bound on \(B_{3R}\).
We use the unit-radius conclusion after rescaling;
metric comparability alone is not the hypothesis
of this estimate.

We give the remaining parabolic step explicitly.
In the fixed holomorphic frames the rescaled metric
equation has a uniformly parabolic scalar principal
part on matrix entries, with coefficients determined
by \(\omega_n\). The remaining terms are smooth
quadratic expressions in the first spatial derivatives,
the positive matrix and its inverse, together with
the smooth zeroth-order term \(2\kappa r_n^2G_n\).
The spatial gradient bound makes these terms uniformly
bounded. Interior parabolic \(W^{2,1}_p\) estimates,
with \(p>6\), give joint
\(C^{1+\beta,(1+\beta)/2}\) bounds on a smaller
cylinder. The right side is then Hölder, and
parabolic Schauder estimates give
\(C^{2+\beta,1+\beta/2}\). Differentiating the smooth
equation and repeating on nested cylinders bounds
every fixed parabolic derivative. At each step there
is positive separation from the spatial and temporal
boundaries. The reference matrix is independent of
\(\tau\), so it contributes no reference velocity.
This is the local bootstrap in
\cite[proof of Theorem~2.5, (2.53)--(2.59)]{XiongYangYau26},
with its prescribed-tensor term replaced here by
the displayed smooth zeroth-order term.

We have proved smooth precompactness on smaller
cylinders. Every subsequential smooth limit is the
already determined uniform limit \(H_{\mathrm{loc}}\).
A subsequence contradiction gives full smooth
convergence at \(\tau=0\) on \(B_R\). Since \(R\)
and the sequence \(T_n\) were arbitrary, this proves
the smooth convergence in Theorem~\ref{thm:energy-curvature}.
No rate for the higher derivatives is needed or asserted.

\subsection{Defect energy and the two radii}

Chern curvature depends smoothly on the metric, its inverse
and its first two derivatives. The smooth inner convergence
therefore implies convergence of curvature in the rescaled
holomorphic frames and fibre metrics. Four-dimensional
curvature energy is unchanged by the constant spatial
dilation and corresponding base rescaling. By
\eqref{eq:local-densities},
\[
 \int_{B_R}|F_{H_{\mathrm{loc}}}|^2\,\dd V_{\mathrm E}
 =48\pi^2\int_0^{R^2}\frac{s}{(1+s)^4}\,\dd s
 =8\pi^2\mathfrak F(R).
\]
This proves \eqref{eq:actual-core-energy}.

For completeness these local limits account for the full
defect charge. An integrable unitary connection on a
Kähler surface has self-dual part \((\Lambda F)\omega/2\).
With our real norm, the self-dual decomposition gives
\begin{equation}\label{eq:chern-weil-energy}
 \int_X|F_H|^2\,\dd V
 =\int_X\tr(F_H\wedge F_H)
       +\int_X|\Lambda F_H|^2\,\dd V .
\end{equation}
The second term is \(2\kappa^2v_0+\|\Phi_H\|_2^2\).
If \(N=\#\Sigma\), the exact sequence implies
\[
 \operatorname{ch}_2(V)
 =\operatorname{ch}_2(L\oplus M)-N[\mathrm{pt}].
\]
Since
\(\operatorname{ch}_2=-\tr(F\wedge F)/(8\pi^2)\),
the topological first terms for \(V\) and \(L\oplus M\)
differ by \(8\pi^2N\). Thus
\begin{equation}\label{eq:total-curvature-limit}
 \lim_{T\to\infty}\int_X|F_{H(T)}|^2\,\dd V
 =\int_X|F_{D_*}|^2\,\dd V+8\pi^2N.
\end{equation}
This also checks the sign: the anti-self-dual local
model contributes positive energy \(8\pi^2\).

Fix disjoint radius-\(r\) coordinate balls and a fixed
\(R>0\). The corrected connection convergence in
Section~\ref{sec:geometry} gives convergence of the
gauge-invariant curvature energy on their fixed
exterior. Subtracting that exterior energy and the
core limits from \eqref{eq:total-curvature-limit}
therefore gives
\begin{align}
 \lim_{T\to\infty}\sum_p
  \int_{B_r(p)\setminus B_{Rr_p(T)}(p)}
             |F_{H(T)}|^2\,\dd V
 ={}&8\pi^2N\{1-\mathfrak F(R)\}\notag\\
    &+\sum_p\int_{B_r(p)}|F_{D_*}|^2\,\dd V .
                                          \label{eq:ordered-tightness}
\end{align}
Let \(r\downarrow0\) and then \(R\uparrow\infty\).
The right side tends to zero, proving tightness in
this ordered-limit sense. Nonnegativity gives the
same conclusion at each point. We assert no
quantitative tail for a simultaneously varying \(R\).

Finally \(\mathfrak F\) is continuous and strictly
increasing on \(R>0\), and \(\mathfrak F(1)=1/2\).
For any fixed \(0<\delta<1\), the actual energies
inside \((1-\delta)r_p(T)\) and \((1+\delta)r_p(T)\)
are eventually below and above \(4\pi^2\), respectively.
Energy in a ball is continuous and nondecreasing
in its radius, so the first radius reaching that
value lies between the two. Letting
\(\delta\downarrow0\) proves \eqref{eq:half-charge-radius}
and completes the proof of
Theorem~\ref{thm:energy-curvature}.
By Corollary~\ref{cor:c0-inner}, the first half-height
inclusion radius has the same leading value, uniformly
in direction. The energy threshold here is half
the defect quantum, not half the total curvature
energy of the compact bundle.
